\section{全ての Poisson 方向の代数的圏量子化}
\label{sec:tilting-quantization}

この節では、滑らかで固有な dg 圏の変形という意味での非可換代数多様体を扱う。
半古典極限は
\[
 HH^2(X)\simeq
 H^0(X,\wedge^2T_X)\oplus H^1(X,T_X)\oplus H^2(X,\mathcal O_X)
\]
による Kodaira--Spencer 類で指定する。
以下の二つの多様体では後の二項は零である。
標識した古典的ファイバーの変形底を用い、自己同値による大域的な商は別に考える。

\subsection{有限次元の tilting 代数}

$X=\mathbb P^1\times\mathbb P^2$ の場合には
\[
 \mathcal E_X=\bigoplus_{\substack{0\leq a\leq1\\0\leq b\leq2}}\mathcal O_X(a,b),
 \qquad
 A_X=\operatorname{End}_X(\mathcal E_X)
\]
とし、$X=(\mathbb P^1)^3$ の場合には
\[
 \mathcal E_X=\bigoplus_{a,b,c\in\{0,1\}}\mathcal O_X(a,b,c),
 \qquad
 A_X=\operatorname{End}_X(\mathcal E_X)
\]
とする。接束を $T_X$ と書く。
射影空間上の標準的な full strong exceptional collection の外積を用いているので、
\[
 \operatorname{Perf}(X)\simeq\operatorname{Perf}(A_X).
\]
$B_n=\operatorname{End}_{\mathbb P^n}
(\bigoplus_{i=0}^{n}\mathcal O(i))$ と置けば、それぞれ
$A_X=B_1\otimes B_2$、$A_X=B_1^{\otimes3}$ である。

直和因子を添字集合 $I$ の元 $\alpha$ で表す。
$I$ は $\{0,1\}\times\{0,1,2\}$ または $\{0,1\}^3$ であり、
座標ごとの半順序を入れる。
\[
 V_{\beta\alpha}
 =\operatorname{Hom}_X(\mathcal O_X(\alpha),\mathcal O_X(\beta))
\]
は $\alpha\leq\beta$ の場合だけ非零である。
$X=\prod_i\mathbb P^{n_i}$ と書くと
\[
 \dim V_{\beta\alpha}
 =\prod_i\binom{n_i+\beta_i-\alpha_i}{n_i}.
\]
対角部分 $E=\bigoplus_{\alpha\in I}\mathbb C e_\alpha$ と、
狭義上三角部分 $J=\bigoplus_{\alpha<\beta}V_{\beta\alpha}$ を固定する。
各 $V_{\beta\alpha}$ に内部次数 $|\beta|-|\alpha|$ を与える。

\begin{lemma}\label{lem:directed-quantization}
$E$ の積、$E$ による左右作用、および上記の corner 空間を保つ任意の結合的積に対し、
$A=E\oplus J$ は有限次元であり、$A$ の大域次元は $3$ 以下である。
相対 normalized bar resolution は長さ $3$ の有限生成射影 $A$-双加群分解となる。
同じ構成は任意の可換基底環上で成り立ち、相対的な smoothness と properness を与える。
\end{lemma}

\begin{proof}
半順序集合 $I$ の狭義鎖の長さは高々 $3$ であるから
$J^{\otimes_E4}=0$ である。
射影 $A\to E$ は代数準同型であり、相対 normalized bar resolution は
\[
 0\longrightarrow A\otimes_EJ^{\otimes_E3}\otimes_EA
 \longrightarrow A\otimes_EJ^{\otimes_E2}\otimes_EA
 \longrightarrow A\otimes_EJ\otimes_EA
 \longrightarrow A\otimes_EA
 \longrightarrow A\longrightarrow0
\]
となる。この複体の完全性は通常の bar contraction による。
各項は $Ae_\beta\otimes_{\mathbb C}e_\alpha A$ の有限直和であるから
$A$-双加群として射影的である。
同様の左加群 bar resolution を任意の $A$-加群に適用すれば大域次元の上界を得る。
可換基底環上でも corner 空間を有限自由加群に置き換えることで同じ証明が使える。
\end{proof}

古典的な積は Koszul である。Koszul 条件は、この有限 bar 複体について
$i\ne j$ に対する $\operatorname{Tor}^{A}_i(E,E)_j=0$ と書ける。
内部次数は高々 $3$ なので有限個の消失条件であり、古典的な積の Zariski 開近傍で保たれる。
Koszul 性のもとでは、corner Hilbert 行列の逆行列から
生成元、二次関係式、三次シジジーの数が決まる。
\[
\begin{array}{c|r|r|r|r|r}
 X&\dim A_X&\text{頂点}&\text{矢}&\text{二次関係式}
      &\text{三次シジジー}\\ \hline
 \mathbb P^1\times\mathbb P^2&60&6&18&18&6\\
 (\mathbb P^1)^3&64&8&24&24&8
\end{array}
\]
最後の列は極小頂点と極大頂点の間の $\operatorname{Tor}_3$ の次元である。
従ってこの開近傍では大域次元はちょうど $3$ となる。
数値は \texttt{check\_tilting.py} で整数演算により再現できる。

\subsection{Jacobi 概型と半普遍変形底}

\[
 C_X=\operatorname{Pois}(X)
 =\{\pi\in H^0(X,\wedge^2T_X)\mid[\pi,\pi]=0\}
\]
を Jacobi 方程式が定義するアフィン概型とする。
この記号にはその非被約構造も含める。

\begin{proposition}\label{prop:poisson-hull}
$A_X$ の結合代数変形の代数的な半普遍底 $(S_X,s_0)$ を取り、
古典的な点での標識を固定できる。
このとき
\[
 \widehat{\mathcal O}_{S_X,s_0}
 \simeq
 \widehat{\mathcal O}_{C_X,0}
\]
という完備局所環の同型があり、接空間上の同型は HKR 同型
$HH^2(A_X)\simeq H^0(X,\wedge^2T_X)$ と整合する。
\end{proposition}

\begin{proof}
$\mathbb P^1$ 上の外積接束は $\mathcal O,\mathcal O(2)$ であり、
$\mathbb P^2$ 上では $\mathcal O,T_{\mathbb P^2},\mathcal O(3)$ である。
Euler 完全系列と K\"unneth の公式から
\[
 H^q(X,\wedge^pT_X)=0\qquad(q>0)
\]
を得る。従って global polyvector field の dg Lie 代数
\[
 \mathfrak g_X^{p-1}=H^0(X,\wedge^pT_X),\qquad d=0
\]
は derived global polyvector field を計算する。
Van den Bergh の大域的 formality 定理
\cite[定理1.1および式(1.3)]{VdBGlobal} と tilting による変形理論の同値
\cite[定理3.18、命題3.21]{BW2024} により、
$\mathfrak g_X$ は $A_X$ の Hochschild dg Lie 代数と $L_\infty$ 準同型の意味で同値である。
$E$ は分離的なので、$E$-relative normalized Hochschild 複体を用いてよい。
実際、形式変形の中で $E$ の原始冪等元は持ち上がり、その持ち上げは
古典的なファイバー上で恒等な共役を除いて一意である。
平坦性によって各 corner の階数も保たれるので、全ての結合代数の形式変形を
この固定した $E$ と corner 空間を持つ表示に移せる。

この同値を、有限型の変形底に結び付ける。
$C_E^n=\operatorname{Hom}_{E\text{-}E}(J^{\otimes_E n},A_X)$ とし、
$\mu_0$ を古典的な積、$d=[\mu_0,-]$ と書く。
$\mu_0+C_E^2$ の中で結合律を満たす積の概型を $M_X$ とする。
これは有限個の二次方程式が定義するアフィン概型である。
$B_E^2=dC_E^1$ の線形補空間 $L\subset C_E^2$ を選び、
\[
 S_X=M_X\cap(\mu_0+L)
\]
と置くと、$T_{s_0}S_X\simeq HH^2(A_X)$ となる。

この切片が hull を与えることを確認する。
corner ごとの基底変換群が $M_X$ に作用し、その無限小作用の像は $B_E^2$ である。
$C_E^1$ 内で $\ker d$ の補空間を選ぶと、この補空間から $B_E^2$ への微分は同型である。
形式的陰関数定理によって、古典的な積の形式近傍で基底変換方向を取り除ける。
小拡大に沿って基底変換を持ち上げ、$B_E^2$ 成分を逐次消せるので、
$\widehat{S_X}_{s_0}$ から結合代数の変形関手への射は形式的に滑らかであり、
接空間上で同型を与える。

一方、$\mathfrak g_X$ の微分は零であるため、その Maurer--Cartan 方程式は
まさに $[\pi,\pi]=0$ である。
gauge 群は小拡大に沿って持ち上がり、原点の一次接空間には自明に作用するので、
$\widehat{C_X}_0$ も対応する変形関手の hull である。
ここで用いる gauge 群は、基底の極大イデアルを $\mathfrak m$ として
$\exp(\mathfrak m\otimes\mathfrak g_X^0)$ であり、標識した古典的ファイバー上で恒等となる。
上記の $L_\infty$ 同値と hull の一意性
\cite[命題2.9]{Schlessinger} によって所要の完備局所環の同型を得る。
\end{proof}

\begin{theorem}\label{thm:algebraic-quantization-all}
$X=\mathbb P^1\times\mathbb P^2$ または $(\mathbb P^1)^3$ とする。
点付き有限型アフィン概型 $(U,u_0)$ と、両方とも \'etale な射
\[
 (C_X,0)\xleftarrow{\ p\ }(U,u_0)
 \xrightarrow{\ q\ }(S_X,s_0)
\]
および $U$ 上の有限自由な結合代数 $\mathcal A$ が存在し、次を満たす。
\begin{enumerate}
 \item $\mathcal A_{u_0}\simeq A_X$ であり、
 $\operatorname{Perf}(\mathcal A_{u_0})\simeq\operatorname{Perf}(X)$ である。
 \item $\mathcal A$ の階数はそれぞれ $60,64$ である。
 $U$ を縮小すると全ての幾何学的ファイバーは Koszul、大域次元 $3$ となり、
 $\operatorname{Perf}(\mathcal A)$ は相対的に smooth かつ proper である。
 \item 古典的な点の Kodaira--Spencer 写像は同型
 \[
 T_{u_0}U\simeq H^0(X,\wedge^2T_X)
       \xrightarrow{\ \sim\ }HH^2(A_X)
 \]
 である。その次元はそれぞれ $34,27$ である。
 \item 任意の $\pi\in C_X(\mathbb C)$ に対し、滑らかな点付き代数曲線
 $(B,b_0)$ 上の $\mathcal A$ の引戻しで、
 局所座標 $t$ に関する一次変形類が $\pi$ となるものが存在する。
\end{enumerate}
\end{theorem}

\begin{proof}
命題 \ref{prop:poisson-hull} の完備局所環の同型に Artin 近似を適用する。
具体的には \cite[補題37.39.2、Tag 0CAT]{StacksArtin} を $N=2$ で用いる。
これにより、与えた同型と一次まで整合する共通の \'etale 近傍 $U$ を得る。
$S_X\subset M_X$ の普遍的な積を $U$ に引き戻して $\mathcal A$ を定める。
その基礎加群は $A_X\otimes_{\mathbb C}\mathcal O_U$ である。
補題 \ref{lem:directed-quantization} とその後の Koszul 開条件が第二項を与える。

HKR 次元は
\[
 h^0(\wedge^2T_{\mathbb P^1\times\mathbb P^2})
 =3\cdot8+10=34,\qquad
 h^0(\wedge^2T_{(\mathbb P^1)^3})=3\cdot3\cdot3=27
\]
である。Artin 近似を二次の極大イデアルまで整合させたので、
Kodaira--Spencer 写像は指定した HKR 同型になる。

最後に直線 $\mathbb A^1\to C_X$, $t\mapsto t\pi$ に沿って $U$ を基底変換する。
点 $(u_0,0)$ の近傍 $B$ は $\mathbb A^1$ 上 \'etale なので滑らかな代数曲線である。
この曲線に沿う一次変形類は $\pi$ である。
\end{proof}

\begin{remark}\label{rem:quantization-scope}
定理は、全ての Poisson 成分を含む、古典的な点の \'etale 局所的な代数的圏量子化を与える。
命題 \ref{prop:poisson-hull} は Jacobi 概型の非被約構造を保持するので、
形式局所的な既約成分と冪零方向も半普遍変形底に引き継がれる。
特定の形式的 star product の全係数との一致を指定する場合には、
その形式族自体の代数化を追加して調べる必要がある。
また、偏極、多重斉次座標環、あるいは指定した connected AS-regular ambient を伴う構成では、
それらの追加構造も変形の対象に含める必要がある。
この節で得た代数の表示は、有限次元の tilting 代数とその有限本の関係式によるものである。
\end{remark}
